\chapter{Reimplementación del Sistema}
\label{cap:reimplementacion}

Este capítulo describe el proceso de reimplementación del módulo TLR, desde las decisiones de diseño hasta la validación de equivalencia funcional. Se presenta de forma racionalizada: no como una narrativa cronológica, sino como el proceso que seguiríamos si tuviéramos que hacerlo de nuevo.

\section{Objetivo de la Reimplementación}

El objetivo no es crear un sistema ``mejor'' que Apollo, sino obtener un \textbf{subsistema equivalente sobre el cual tengamos control total}:

\begin{itemize}
    \item \textbf{Fidelidad:} Producir los mismos resultados que Apollo para las mismas entradas
    \item \textbf{Simplicidad:} Eliminar dependencias innecesarias (CyberRT, Docker, Bazel)
    \item \textbf{Instrumentabilidad:} Poder agregar logging, visualizaciones, modificaciones
    \item \textbf{Testabilidad:} Poder ejecutar y testear fácilmente sin configuración compleja
\end{itemize}

\section{Decisiones de Diseño}

\subsection{Framework: PyTorch}

Se eligió PyTorch sobre Caffe por:

\begin{itemize}
    \item \textbf{Mantenimiento activo:} PyTorch está activamente desarrollado; Caffe está obsoleto
    \item \textbf{Debugging:} Ejecución dinámica permite debugging con herramientas Python estándar
    \item \textbf{Ecosistema:} Mayor disponibilidad de recursos y documentación
\end{itemize}

\subsection{Estructura de Módulos}

\begin{verbatim}
src/tlr/
├── pipeline.py              # Pipeline principal
├── detector.py              # Módulo de detección
├── recognizer.py            # Módulo de reconocimiento
├── tracking.py              # Sistema de tracking
├── selector.py              # Algoritmo húngaro
├── hungarian_optimizer.py   # Implementación de Munkres
├── utils.py                 # Utilidades
├── weights/                 # Modelos pre-entrenados
└── confs/                   # Configuraciones
\end{verbatim}

\subsection{API Simplificada}

\begin{codesnippet}
\begin{verbatim}
from tlr.pipeline import load_pipeline

pipeline = load_pipeline('cuda:0')  # o 'cpu'
result = pipeline(image, projection_boxes)
\end{verbatim}
\end{codesnippet}

Esta API contrasta con Apollo que requiere Docker, CyberRT, y configuraciones extensas.

\section{Proceso de Reimplementación}

\subsection{Etapa 1: Conversión de Modelos}

Los modelos de Apollo están en formato Caffe. El proceso de conversión:

\begin{enumerate}
    \item Parsear archivos \texttt{.prototxt} para extraer arquitectura
    \item Crear capas equivalentes en PyTorch
    \item Cargar pesos desde \texttt{.caffemodel}
    \item Validar equivalencia numérica capa por capa
\end{enumerate}

Se convirtieron 4 modelos: detector (\texttt{tl.torch}), clasificador vertical (\texttt{vert.torch}), horizontal (\texttt{hori.torch}), y quad (\texttt{quad.torch}).

\subsection{Etapa 2: Implementación del Pipeline}

Cada etapa del pipeline se implementó siguiendo exactamente el código de Apollo:

\textbf{Preprocessing:}
\begin{itemize}
    \item Expansión de crops: \texttt{crop\_scale=2.5}, \texttt{min\_crop\_size=270}
    \item Normalización: means BGR [102.98, 115.95, 122.77]
\end{itemize}

\textbf{Detection:}
\begin{itemize}
    \item Red CNN con output [scores + coordenadas]
    \item NMS en tres niveles (RPN, RCNN, Global) con threshold 0.6
\end{itemize}

\textbf{Assignment:}
\begin{itemize}
    \item Algoritmo húngaro con pesos 70\% distancia, 30\% confianza
    \item Distancia gaussiana 2D ($\sigma=100$)
    \item Clipping de scores a máximo 0.9
\end{itemize}

\textbf{Recognition:}
\begin{itemize}
    \item Selección de clasificador según orientación detectada
    \item Normalización específica: means BGR [66.56, 66.58, 69.06], scale 0.01
    \item Mapeo de probabilidades a color con threshold 0.5
\end{itemize}

\textbf{Tracking:}
\begin{itemize}
    \item Historial por signal\_id (diccionario Python)
    \item Implementación de las 4 reglas de seguridad
\end{itemize}

\subsection{Etapa 3: Simplificación de Interfaces}

\textbf{Projection boxes:} En lugar de obtenerlas dinámicamente del HD-Map, se leen de un archivo de texto pre-calculado. Esto no afecta la lógica core: desde el punto de vista del TLR, el input es el mismo (lista de bounding boxes).

\textbf{Formato del archivo:}
\begin{codesnippet}
\begin{verbatim}
# Una línea por frame: x1 y1 x2 y2 [x1 y1 x2 y2 ...]
100 50 150 120 500 80 550 150
102 52 152 122 498 82 548 152
\end{verbatim}
\end{codesnippet}

\section{Validación de Equivalencia}

La validación de equivalencia se realizó en dos niveles:

\subsection{Validación por Construcción}

La reimplementación se hizo traduciendo directamente del código C++ de Apollo:

\begin{itemize}
    \item Mismo orden de operaciones
    \item Mismos valores numéricos exactos
    \item Mismas reglas de negocio
    \item Comentarios en código referenciando líneas de Apollo
\end{itemize}

\subsection{Validación mediante Tests de Corrección}

Durante el desarrollo, ejecutamos el sistema con inputs reales y corregimos discrepancias. Este proceso reveló 4 errores que fueron corregidos:

\subsubsection{Error 1: Detección de Blink}

\textbf{Problema:} El blink (parpadeo) se implementó inicialmente para el amarillo.

\textbf{Corrección:} El blink se detecta solo en verde. Revisión del código de Apollo confirmó esta lógica.

\subsubsection{Error 2: Secuencia de Colores}

\textbf{Problema:} Asumimos secuencia RED→YELLOW→GREEN.

\textbf{Corrección:} La secuencia correcta es GREEN→YELLOW→RED. El amarillo solo aparece después del verde, nunca después del rojo.

\subsubsection{Error 3: Pesos del Húngaro}

\textbf{Problema:} Los pesos estaban invertidos (30\% distancia, 70\% confianza).

\textbf{Corrección:} Deben ser 70\% distancia, 30\% confianza.

\subsubsection{Error 4: Orden de Means BGR}

\textbf{Problema:} Los means de normalización estaban en orden RGB en lugar de BGR.

\textbf{Corrección:} OpenCV devuelve imágenes en BGR. Los means deben aplicarse en ese orden.

\subsection{Impacto de las Correcciones}

Después de corregir estos errores:

\begin{itemize}
    \item El recognizer clasifica correctamente en condiciones normales
    \item Las asignaciones del húngaro son consistentes con la proximidad espacial
    \item Las reglas de tracking se aplican en los momentos correctos
    \item El blink detection funciona cuando se prueba con semáforos parpadeando en verde
\end{itemize}

\section{Diferencias Estructurales}

Las únicas diferencias con Apollo son estructurales, no funcionales:

\begin{table}[H]
\centering
\begin{tabular}{|l|l|l|}
\hline
\textbf{Aspecto} & \textbf{Apollo} & \textbf{Reimplementación} \\
\hline
Lenguaje & C++ & Python \\
Framework & Caffe & PyTorch \\
Tracking storage & Vector (O(n)) & Diccionario (O(1)) \\
NMS ordering & Ascendente desde atrás & Descendente desde adelante \\
Projection boxes & HD-Map dinámico & Pre-calculadas \\
\hline
\end{tabular}
\caption{Diferencias estructurales (no funcionales)}
\end{table}

Todas estas diferencias son matemáticamente equivalentes.

\section{Comparación de Pipelines}

Las siguientes figuras muestran la estructura de ambos pipelines y su relación.

% Diagramas de comparacion de pipelines
% Para usar en el capitulo 5 de la tesis (Aislamiento del Modulo TLR)

% IMPORTANTE: Desactivar shorthand de babel para que TikZ funcione
\shorthandoff{<>}

% ============================================================
% DIAGRAMA 1: PIPELINE APOLLO ORIGINAL
% ============================================================

\begin{figure}[H]
\centering
\begin{tikzpicture}[
    node distance=0.8cm,
    box/.style={rectangle, draw, rounded corners, minimum width=3cm, minimum height=0.8cm, align=center, font=\small},
    inputbox/.style={box, fill=blue!15},
    procbox/.style={box, fill=orange!20},
    outputbox/.style={box, fill=green!15},
    arrow/.style={->, thick, >=stealth}
]

% Inputs
\node[inputbox] (img) {Imagen de Camara};
\node[inputbox, right=0.5cm of img] (hdmap) {HD-Map\\(Coords 3D)};
\node[inputbox, right=0.5cm of hdmap] (pose) {Pose Vehiculo\\(GPS + TF)};
\node[inputbox, right=0.5cm of pose] (calib) {Calibracion\\Camara};

% Etapas
\node[procbox, below=1.8cm of hdmap, xshift=0.5cm] (prep) {\textbf{Region Proposal}\\Query HD-Map, Proyeccion 3D a 2D};

\node[procbox, below=0.8cm of prep] (detect) {\textbf{Detection}\\CNN SSD-style + NMS};

\node[procbox, below=0.8cm of detect] (select) {\textbf{Assignment}\\Algoritmo Hungaro\\(70\% dist, 30\% conf)};

\node[procbox, below=0.8cm of select] (recog) {\textbf{Recognition}\\3 Clasificadores CNN\\(vert/quad/hori)};

\node[procbox, below=0.8cm of recog] (track) {\textbf{Tracking}\\Reglas temporales\\+ Semantic Voting};

% Output
\node[outputbox, below=0.8cm of track] (output) {\textbf{Output}\\Color + Blink + Signal ID};

% Coordenada para la linea horizontal comun
\coordinate (linea_comun) at ($(img.south)+(0,-0.5)$);

% Arrows inputs - todas bajan a la misma altura y luego van al centro
\draw[arrow] (img.south) -- (img.south |- linea_comun) -| (prep.north);
\draw[arrow] (hdmap.south) -- (hdmap.south |- linea_comun) -- (prep.north |- linea_comun) -- (prep.north);
\draw[arrow] (pose.south) -- (pose.south |- linea_comun) -| (prep.north);
\draw[arrow] (calib.south) -- (calib.south |- linea_comun) -| (prep.north);

% Arrows between stages
\draw[arrow] (prep) -- (detect);
\draw[arrow] (detect) -- (select);
\draw[arrow] (select) -- (recog);
\draw[arrow] (recog) -- (track);
\draw[arrow] (track) -- (output);

% Modulos Apollo
\node[right=0.3cm of prep, font=\tiny\ttfamily, text=gray] {traffic\_light\_region\_proposal};
\node[right=0.3cm of detect, font=\tiny\ttfamily, text=gray] {traffic\_light\_detection};
\node[right=0.3cm of select, font=\tiny\ttfamily, text=gray] {traffic\_light\_detection};
\node[right=0.3cm of recog, font=\tiny\ttfamily, text=gray] {traffic\_light\_recognition};
\node[right=0.3cm of track, font=\tiny\ttfamily, text=gray] {traffic\_light\_tracking};

\end{tikzpicture}
\caption{Pipeline de Apollo original con los 4 modulos}
\label{fig:pipeline_apollo}
\end{figure}


% ============================================================
% DIAGRAMA 2: PIPELINE REIMPLEMENTACION
% ============================================================

\begin{figure}[H]
\centering
\begin{tikzpicture}[
    node distance=0.8cm,
    box/.style={rectangle, draw, rounded corners, minimum width=3cm, minimum height=0.8cm, align=center, font=\small},
    inputbox/.style={box, fill=blue!15},
    procbox/.style={box, fill=orange!20},
    outputbox/.style={box, fill=green!15},
    simplebox/.style={box, fill=yellow!20, dashed},
    arrow/.style={->, thick, >=stealth}
]

% Inputs (simplificados)
\node[inputbox] (img) {Imagen de Camara};
\node[simplebox, right=1cm of img] (boxes) {Projection Boxes\\(archivo .txt)};

% Etapas
\node[procbox, below=1cm of img, xshift=1.5cm] (prep) {\textbf{Preprocessing}\\Crop + Resize 270x270};

\node[procbox, below=0.8cm of prep] (detect) {\textbf{Detection}\\CNN SSD-style + NMS};

\node[procbox, below=0.8cm of detect] (select) {\textbf{Assignment}\\Algoritmo Hungaro\\(70\% dist, 30\% conf)};

\node[procbox, below=0.8cm of select] (recog) {\textbf{Recognition}\\3 Clasificadores CNN\\(vert/quad/hori)};

\node[procbox, below=0.8cm of recog] (track) {\textbf{Tracking}\\Reglas temporales};

% Output
\node[outputbox, below=0.8cm of track] (output) {\textbf{Output}\\Color + Blink + Signal ID};

% Arrows inputs
\draw[arrow] (img.south) -- ++(0,-0.3) -| (prep.north);
\draw[arrow] (boxes.south) -- ++(0,-0.3) -| (prep.north);

% Arrows between stages
\draw[arrow] (prep) -- (detect);
\draw[arrow] (detect) -- (select);
\draw[arrow] (select) -- (recog);
\draw[arrow] (recog) -- (track);
\draw[arrow] (track) -- (output);

% Archivos Python
\node[right=0.3cm of prep, font=\tiny\ttfamily, text=gray] {utils.py};
\node[right=0.3cm of detect, font=\tiny\ttfamily, text=gray] {detector.py};
\node[right=0.3cm of select, font=\tiny\ttfamily, text=gray] {selector.py};
\node[right=0.3cm of recog, font=\tiny\ttfamily, text=gray] {recognizer.py};
\node[right=0.3cm of track, font=\tiny\ttfamily, text=gray] {tracking.py};

% Nota de simplificacion
\node[below=0.3cm of output, font=\footnotesize, text=gray, align=center] {Linea punteada = simplificacion\\respecto a Apollo};

\end{tikzpicture}
\caption{Pipeline del modulo aislado (con interfaz simplificada)}
\label{fig:pipeline_aislado}
\end{figure}


% ============================================================
% DIAGRAMA 3: COMPARACION LADO A LADO
% ============================================================

\begin{figure}[H]
\centering
\begin{tikzpicture}[
    node distance=0.6cm,
    box/.style={rectangle, draw, rounded corners, minimum width=2.2cm, minimum height=0.6cm, align=center, font=\scriptsize},
    inputbox/.style={box, fill=blue!15},
    procbox/.style={box, fill=orange!20},
    outputbox/.style={box, fill=green!15},
    simplebox/.style={box, fill=yellow!20, dashed},
    greybox/.style={box, fill=gray!30},
    arrow/.style={->, thick, >=stealth},
    eqarrow/.style={<->, thick, >=stealth, color=blue!60}
]

% ========== APOLLO (izquierda) ==========
\node[font=\bfseries] (titleA) {Apollo Original};

\node[inputbox, below=0.3cm of titleA, xshift=-2cm] (imgA) {Imagen};
\node[inputbox, right=0.5cm of imgA] (hdmapA) {HD-Map};
\node[inputbox, right=0.5cm of hdmapA] (poseA) {Pose};

\node[procbox, below=1.2cm of hdmapA] (prepA) {Region Proposal\\(3D a 2D)};
\node[procbox, below=0.5cm of prepA] (detectA) {Detection};
\node[procbox, below=0.5cm of detectA] (selectA) {Assignment};
\node[procbox, below=0.5cm of selectA] (recogA) {Recognition};
\node[procbox, below=0.5cm of recogA] (trackA) {Tracking};
\node[outputbox, below=0.5cm of trackA] (outA) {Output};

% Coordenada para linea horizontal Apollo
\coordinate (linea_A) at ($(imgA.south)+(0,-0.3)$);

% Arrows Apollo - todas a la misma altura
\draw[arrow] (imgA.south) -- (imgA.south |- linea_A) -| (prepA.north);
\draw[arrow] (hdmapA.south) -- (hdmapA.south |- linea_A) -- (prepA.north |- linea_A) -- (prepA.north);
\draw[arrow] (poseA.south) -- (poseA.south |- linea_A) -| (prepA.north);
\draw[arrow] (prepA) -- (detectA);
\draw[arrow] (detectA) -- (selectA);
\draw[arrow] (selectA) -- (recogA);
\draw[arrow] (recogA) -- (trackA);
\draw[arrow] (trackA) -- (outA);

% ========== REIMPLEMENTACION (derecha) ==========
\node[font=\bfseries, right=5cm of titleA] (titleR) {Modulo Aislado};

\node[inputbox, below=0.3cm of titleR, xshift=-0.8cm] (imgR) {Imagen};
\node[simplebox, right=0.5cm of imgR] (boxesR) {Boxes (txt)};

\node[simplebox, below=1.2cm of imgR, xshift=0.8cm] (prepR) {Preprocessing\\(2D directo)};
\node[procbox, below=0.5cm of prepR] (detectR) {Detection};
\node[procbox, below=0.5cm of detectR] (selectR) {Assignment};
\node[procbox, below=0.5cm of selectR] (recogR) {Recognition};
\node[procbox, below=0.5cm of recogR] (trackR) {Tracking};
\node[outputbox, below=0.5cm of trackR] (outR) {Output};

% Coordenada para linea horizontal Reimpl
\coordinate (linea_R) at ($(imgR.south)+(0,-0.3)$);

% Arrows Reimpl - todas a la misma altura
\draw[arrow] (imgR.south) -- (imgR.south |- linea_R) -| (prepR.north);
\draw[arrow] (boxesR.south) -- (boxesR.south |- linea_R) -| (prepR.north);
\draw[arrow] (prepR) -- (detectR);
\draw[arrow] (detectR) -- (selectR);
\draw[arrow] (selectR) -- (recogR);
\draw[arrow] (recogR) -- (trackR);
\draw[arrow] (trackR) -- (outR);

% ========== EQUIVALENCIAS ==========
\draw[eqarrow] (detectA.east) -- node[above, font=\tiny] {=} (detectR.west);
\draw[eqarrow] (selectA.east) -- node[above, font=\tiny] {=} (selectR.west);
\draw[eqarrow] (recogA.east) -- node[above, font=\tiny] {=} (recogR.west);
\draw[eqarrow] (trackA.east) -- node[above, font=\tiny] {=} (trackR.west);

% Diferencia
\draw[dashed, red, thick] (prepA.east) -- node[above, font=\tiny, red] {$\neq$} (prepR.west);

% Leyenda
\node[below=0.8cm of outA, xshift=2cm, font=\scriptsize, align=left] {
    \textcolor{blue}{$\leftrightarrow$ Equivalente funcionalmente}\\
    \textcolor{red}{- - - Simplificado (no HD-Map)}
};

\end{tikzpicture}
\caption[Comparación lado a lado: Apollo vs módulo aislado.]{Comparacion lado a lado: Apollo original vs modulo aislado. Las flechas azules indican etapas funcionalmente equivalentes; la linea roja punteada indica la unica etapa cuya interfaz se simplifico (HD-Map dinamico reemplazado por projection boxes precalculadas).}
\label{fig:pipeline_comparison}
\end{figure}

% Reactivar shorthand de babel
\shorthandon{<>}

\section{Resultado}

El resultado es un sistema que:

\begin{itemize}
    \item \textbf{Es funcionalmente equivalente} a Apollo para la lógica core del TLR
    \item \textbf{Es simple de usar:} \texttt{pip install requirements.txt}, sin Docker ni configuración compleja
    \item \textbf{Es modificable:} Código Python limpio y documentado
    \item \textbf{Es instrumentable:} Fácil agregar logging, visualizaciones, breakpoints
\end{itemize}

Esta reimplementación sirvió como base para el proceso de testing descrito en el siguiente capítulo.